\subsection{Injectivity of the ground-space map}

\begin{theorem}[Ground-space map injectivity]\label{thm:ground_space_map_injective}
    \lean{MPSTensor.groundSpaceMap_injective}
    \leanok
    \uses{def:ground_space_map, def:injective}
    If $A$ is injective and $L \ge 1$, then $\Gamma_L$ is injective.
\end{theorem}

\begin{proof}
    \leanok
    If $\Gamma_L(X) = 0$, then $\tr(A^\sigma X) = 0$ for every length-$L$ word
    $\sigma$. Since $A$ is injective, the products $\{A^\sigma\}_\sigma$ span
    $\MN{D}$, so the trace pairing gives $X = 0$.
\end{proof}

\begin{theorem}[Ground-space map injectivity from word span]
    \label{thm:ground_space_map_injective_word_span}
    \lean{MPSTensor.groundSpaceMap_injective_of_wordSpan_eq_top}
    \leanok
    \uses{def:ground_space_map}
    If the length-$L$ products $\{A^\sigma : |\sigma| = L\}$ span $\MN{D}$, then
    $\Gamma_L$ is injective.
\end{theorem}

\begin{proof}
    \leanok
    If $\Gamma_L(X) = 0$, then $\tr(A^\sigma X) = 0$ for every word $\sigma$ of
    length $L$. The spanning hypothesis turns this into $\tr(MX) = 0$ for every
    $M \in \MN{D}$, and the trace pairing gives $X = 0$.
\end{proof}

\begin{lemma}[Ground-space map injectivity after blocking]
    \label{lem:ground_space_map_injective_blk}
    \lean{MPSTensor.groundSpaceMap_injective_of_isNBlkInjective}
    \leanok
    \uses{thm:ground_space_map_injective_word_span, def:l_blk_injective}
    If the length-$L_0$ products span the full matrix algebra,
    $\spn\{A^\omega : |\omega| = L_0\} = \MN{D}$, then $\Gamma_{L_0}$ is
    injective.
\end{lemma}

\begin{proof}
    \leanok
    The hypothesis is precisely the fixed-length spanning condition needed by
    Theorem~\ref{thm:ground_space_map_injective_word_span}.
\end{proof}

\begin{theorem}[Ground-space dimension]\label{thm:ground_space_finrank_eq}
    \lean{MPSTensor.groundSpace_finrank_eq}
    \leanok
    \uses{thm:ground_space_map_injective, lem:ground_space_finrank_le}
    If $A$ is injective and $L \ge 1$, then $\dim G_L(A) = D^2$.
\end{theorem}

\begin{proof}
    \leanok
    The upper bound $\dim G_L(A) \le D^2$ is
    Lemma~\ref{lem:ground_space_finrank_le}. Injectivity of $\Gamma_L$ gives
    $\dim G_L(A) = \dim \MN{D} = D^2$.
\end{proof}

\subsection{The intersection property}

\begin{theorem}[One inclusion of the intersection property]\label{thm:ground_space_intersection}
    \lean{MPSTensor.groundSpace_intersection}
    \leanok
    \uses{thm:ground_space_map_injective, def:in_left_ground, def:in_right_ground,
        def:ground_space_parent, def:injective}
    If $A$ is injective and $L \ge 2$, then the nontrivial inclusion in the
    intersection property holds:
    \begin{align}
        G_L^{\mathrm{left}}\cap G_L^{\mathrm{right}}
        &\subseteq G_{L+1}(A).
        \notag
    \end{align}
\end{theorem}

\begin{proof}
    \leanok
    The ``inverting and growing back'' argument compares the left and right
    restrictions on their common $(L-1)$-site overlap and then reconstructs the
    full $(L+1)$-site word from a single boundary matrix:
    \begin{center}
        % Formula/source: Perez-Garcia et al., Eq. (eq:injectivity), define
        % Gamma_{L+1}(X') by the coefficients
        % tr(X' A^{sigma_1} ... A^{sigma_{L+1}}).
        % In indices, this is the sum of
        % X'_{a_0 a_1}(A^{sigma_1})_{a_1 a_2} ...
        % (A^{sigma_{L+1}})_{a_{L+1} a_0} over every virtual index a_r.
        % The physical indices sigma_1,...,sigma_{L+1} remain free.
        % Orientation: placing X' first matches the source's cyclic order;
        % the periodic return closes the final A factor back into X'.
        % Semantic boundary signature: (virtual-west=0 | virtual-east=0 |
        % physical-up=L+1 | physical-down=0).  This finite schematic emits
        % (virtual-west=0 | virtual-east=0 | physical-up=3 | physical-down=0).
        % The traced sides draw the periodic return; the site-count brace does
        % not draw in the kernel yet (#5482), so its mathematics stays as the
        % south label under the return.
        \tenkzkernel
        \begin{tenkz}[rows={wire}, cols=5, boundary=periodic]
            \tn[skin=ring]{X'} &
            \tn[label pos=s, ports={90:physical:$\sigma_1$}]{A} &
            \tn[label pos=s, ports={90:physical}]{A} & \tn[skin=dots]{} &
            \tn[label pos=s, ports={90:physical:$\sigma_{L+1}$}]{A}
            \tnmark[form=label, label pos=s]{s outside}{$L+1\text{ sites}$}
        \end{tenkz}
    \end{center}
    \begin{enumerate}
        \item From the right ground condition, for each physical index~$i$ there
              exists a unique $Y_i \in \MN{D}$, by injectivity of $\Gamma_L$,
              such that $\psi(i,\sigma) = \tr(A^\sigma Y_i)$.
        \item Similarly, from the left ground condition, for each~$j$ there
              exists a unique $Z_j$ such that
              $\psi(\sigma,j) = \tr(A^\sigma Z_j)$.
        \item Matching on the $(L-1)$-site overlap and using nondegeneracy of the
              trace pairing gives $A^jY_i = Z_jA^i$ for all $i,j$.
        \item Since $A$ is injective, the matrices $\{A^j\}$ span $\MN{D}$
              (Definition~\ref{def:injective}), so
              $\Id = \sum_j c_jA^j$ for some coefficients $c_j \in \C$. Thus,
              with $X' := \sum_j c_jZ_j$,
              $Y_i = \Id Y_i = \sum_j c_jA^jY_i = \sum_j c_jZ_jA^i = X'A^i$.
        \item By trace cyclicity, $\psi(\sigma) = \tr(A^\sigma X')$, so
              $\psi \in G_{L+1}(A)$.
    \end{enumerate}
\end{proof}

\begin{theorem}[Intersection property]\label{thm:ground_space_iff}
    \lean{MPSTensor.groundSpace_iff_left_right}
    \leanok
    \uses{lem:ground_space_in_left_ground, lem:ground_space_in_right_ground,
        thm:ground_space_intersection}
    For injective $A$ and $L \ge 2$,
    $G_L^{\mathrm{left}}\cap G_L^{\mathrm{right}} = G_{L+1}(A)$,
    equivalently, $\psi \in G_{L+1}(A)$ if and only if
    $\psi \in G_L^{\mathrm{left}}\cap G_L^{\mathrm{right}}$.
\end{theorem}

\begin{proof}
    \leanok
    The claim follows from the two forward-direction lemmas and
    Theorem~\ref{thm:ground_space_intersection}.
\end{proof}

\subsection{Unique ground state}

\subsubsection{Contiguous and cyclic interval restrictions}

\begin{remark}[Contiguous interval restrictions]\label{rem:contiguous_window_operators}
    \lean{MPSTensor.contiguousCfg}
    \lean{MPSTensor.contiguousRestrictₗ}
    \lean{MPSTensor.contiguousRestrictₗ_apply}
    \lean{MPSTensor.contiguousCfg_zero_full}
    \lean{MPSTensor.contiguousRestrictₗ_restrictLast}
    \lean{MPSTensor.contiguousRestrictₗ_restrictFirst}
    \lean{MPSTensor.tailRestrictₗ_contiguousRestrictₗ}
    \leanok
    For a non-wrapping interval $[s,s+M-1] \subseteq \{0,\ldots,N-1\}$ and a
    full configuration $\rho$ supplying the complementary values, the
    restriction map is
    \begin{align}
        (\Res^\rho_{s,M}\psi)(\omega)
        &= \psi\!\left(\rho^{[s,s+M-1]\leftarrow\omega}\right).
        \notag
    \end{align}
    The listed identities are the equations for $s=0$, $M=N$, for deleting the
    first or last site, and for first fixing $K$ sites and then restricting to
    $[s+K,s+K+L-1]$.
\end{remark}

\begin{remark}[Cyclic interval restrictions]\label{rem:cyclic_window_operators}
    \lean{MPSTensor.cyclicCfg}
    \lean{MPSTensor.cyclicRestrictₗ}
    \lean{MPSTensor.cyclicRestrictₗ_apply}
    \lean{MPSTensor.cyclicRestrictₗ_congr_outside}
    \lean{MPSTensor.cyclicRestrictₗ_cyclicCfg_outside}
    \lean{MPSTensor.cyclicCfg_cyclicForwardSite}
    \lean{MPSTensor.eq_cyclic_site_of_offset_eq}
    \lean{MPSTensor.cyclicRestrictₗ_restrictLast}
    \lean{MPSTensor.cyclicRestrictₗ_restrictFirst}
    \lean{MPSTensor.cyclicCfg_eq_contiguousCfg}
    \lean{MPSTensor.cyclicRestrictₗ_eq_contiguousRestrictₗ}
    \leanok
    On the periodic chain, for a full configuration $\rho$ supplying the
    complementary values, put
    \begin{align}
        \delta_i(k)
        &:= (k+N-i)\bmod N,
        \notag\\
        \rho^{\mathrm{cyc}}_{i,L,\omega}(k)
        &:=
        \begin{cases}
            \omega(\delta_i(k)), & \delta_i(k)<L,\\
            \rho(k),             & \delta_i(k)\ge L.
        \end{cases}
        \notag
    \end{align}
    The restriction to the cyclic interval beginning at $i$ is
    \begin{align}
        (\Res^\rho_{i,L}\psi)(\omega)
        &= \psi\!\left(\rho^{\mathrm{cyc}}_{i,L,\omega}\right),
        \notag
    \end{align}
    with all site labels read modulo $N$. If $L\le N$, this is the same
    configuration as $\rho^{[i,i+L-1]\leftarrow\omega}$. If two boundary
    conditions agree whenever $\delta_i(k)\ge L$, their restrictions are equal.
    In particular, filling the interval inside the outside configuration leaves
    the restriction unchanged:
    \begin{align}
        \Res^{\rho^{\mathrm{cyc}}_{i,L,\omega}}_{i,L}\psi
        &= \Res^\rho_{i,L}\psi.
        \notag
    \end{align}
    For every $r<L$, the inserted word is recovered as
    $\rho^{\mathrm{cyc}}_{i,L,\omega}(i+r)=\omega(r)$. Here site labels are
    taken modulo $N$. The remaining listed identities say that deleting the
    final site of an $(L+1)$-interval gives the $L$-interval beginning at $i$;
    if $L+1\le N$, deleting the first site gives the $L$-interval beginning at
    $i+1$; and the cyclic formula agrees with the contiguous formula when the
    interval does not wrap around the ring.
\end{remark}

\begin{lemma}[One-site restrictions of cyclic-window boundary conditions]
    \label{lem:cyclic_window_boundary_matrix_transport}
    \lean{MPSTensor.cyclicRestrictₗ_restrictFirst_groundSpaceMap}
    \lean{MPSTensor.cyclicRestrictₗ_restrictLast_groundSpaceMap}
    \lean{MPSTensor.adjacent_cyclicRestrictₗ_witness_overlap}
    \lean{MPSTensor.adjacent_cyclicRestrictₗ_witness_overlap_common_background}
    \lean{MPSTensor.boundary_witness_product_of_adjacent_overlaps}
    \lean{MPSTensor.adjacent_cyclicRestrictₗ_witness_product}
    \lean{MPSTensor.adjacent_cyclicRestrictₗ_witness_product_common_background}
    \lean{MPSTensor.adjacent_cyclicRestrictₗ_witness_product_common_background_named}
    \leanok
    \uses{rem:cyclic_window_operators, def:ground_space_parent, def:l_blk_injective}
    If a cyclic $(L+1)$-window is represented by a boundary matrix $Y$, then
    deleting one endpoint gives
    \begin{align}
        \Gamma_{L+1}(Y)\big|_{\mathrm{last}=b}
        &= \Gamma_L(A^bY),
        \notag\\
        \Gamma_{L+1}(Y)\big|_{\mathrm{first}=a}
        &= \Gamma_L(YA^a).
        \notag
    \end{align}
    Hence two adjacent cyclic windows whose restrictions agree on the common
    length-$L$ overlap have boundary matrices satisfying $Y_1A^a=A^bY_2$.
    Consequently, for a finite chain of adjacent overlaps indexed by
    $0\le r<n$, satisfying $Y_rA^{a_r}=A^{b_r}Y_{r+1}$, one has
    $Y_0A^{a_0}\cdots A^{a_{n-1}} = A^{b_0}\cdots A^{b_{n-1}}Y_n$.
    In particular, if the two endpoint words are named by the equations
    $a_r=\rho(i_0+r)$ and $b_r=\rho(i_0+r+L+1)$, with site labels read modulo
    $N$, then the same transport identity is written as $Y_0A^a=A^bY_n$.
\end{lemma}

\begin{proof}
    \leanok
    The first two identities are the cyclic first- and last-site deletion
    identities, followed by the corresponding ground-space restriction
    formulas. For adjacent windows, equality of the completed outside
    configurations gives
    $\Gamma_L(Y_0A^a)=\Gamma_L(A^bY_n)$. Injectivity of $\Gamma_L$ gives
    $Y_0A^a=A^bY_n$. Iterating the adjacent identity gives the word-product
    identity.
\end{proof}

\begin{lemma}[Iterated contiguous-window intersection]\label{lem:contiguous_mem_ground_space}
    \lean{MPSTensor.contiguous_mem_groundSpace}
    \leanok
    \uses{thm:ground_space_intersection}
    If an $N$-site state satisfies the $L$-site ground condition on every
    non-wrapping contiguous window, with $L \ge 2$, then it lies in $G_N(A)$.
\end{lemma}

\begin{proof}
    \leanok
    The induction repeatedly merges two neighbouring admissible windows
    $[i,i+L-1]$ and $[i+1,i+L]$ into the larger window $[i,i+L]$ by applying
    Theorem~\ref{thm:ground_space_intersection} to their common overlap
    $[i+1,i+L-1]$.
\end{proof}

\begin{definition}[Suffix restriction]\label{def:tail_restrict_parent}
    \lean{MPSTensor.tailRestrictₗ}
    \leanok
    Fix a prefix $u\in\{0,\ldots,d-1\}^K$. The suffix restriction of an
    $(K+L)$-site state $\psi$ is the $L$-site state given by
    $(R^{\mathrm{tail}}_u\psi)(\sigma)=\psi(u,\sigma)$.
\end{definition}

\begin{definition}[Suffix ground membership]\label{def:in_tail_ground}
    \lean{MPSTensor.InTailGround}
    \leanok
    \uses{def:tail_restrict_parent, def:ground_space_parent}
    A state $\psi$ on $K+L$ sites lies in the suffix ground condition if every
    fixed prefix $u$ gives a suffix state in $G_L(A)$.
\end{definition}

\begin{lemma}[Ground-space suffix restriction]\label{lem:ground_space_in_tail_ground}
    \lean{MPSTensor.restrictLast_groundSpaceMap,
        MPSTensor.restrictFirst_groundSpaceMap,
        MPSTensor.tailRestrictₗ_groundSpaceMap,
        MPSTensor.groundSpace_inTailGround}
    \leanok
    \uses{def:in_tail_ground, def:ground_space_parent}
    If $\psi \in G_{K+L}(A)$, then every fixed-prefix suffix restriction of
    $\psi$ lies in $G_L(A)$. For a vector in $G_{L+1}(A)$, fixing either
    endpoint gives the expected left or right multiplication of its boundary
    matrix.
\end{lemma}

\begin{proof}
    \leanok
    Write $\psi(\rho) = \tr(A^\rho X)$ for a boundary matrix $X$. Fixing the
    final site gives the boundary matrix $A^jX$, while fixing the first site
    gives $XA^i$ by cyclicity of the trace. Fixing a prefix $u$ gives
    $\psi(u,\sigma)=\tr(A^uA^\sigma X)=\tr(A^\sigma XA^u)$, so the suffix
    restriction again has the form of a ground-space vector.
\end{proof}

\begin{lemma}[Left-tail compatibility]\label{lem:left_tail_compatibility}
    \lean{MPSTensor.exists_left_tail_compatibility}
    \leanok
    \uses{def:in_left_ground, def:in_tail_ground, def:ground_space_parent}
    Assume $\psi$ on $K+L_0+1$ sites satisfies the left ground condition on the
    first $K+L_0$ sites and the suffix ground condition on every fixed prefix of
    length $K$. Then there exist matrices $(Z_j)_j$ and $(Y_u)_u$ with the two
    trace representations below and, for all $j$ and $u$, the compatibility
    relation
    \begin{align}
        \psi(\sigma,j)
        &= \tr(A^\sigma Z_j),
        \notag\\
        \psi(u,\tau)
        &= \tr(A^\tau Y_u),
        \notag\\
        Z_jA^u
        &= A^jY_u.
        \notag
    \end{align}
\end{lemma}

\begin{proof}
    \leanok
    Choose $Z_j$ from the long left-window condition and $Y_u$ from the suffix
    ground condition. The two restrictions give the same vector in
    $G_{L_0}(A)$, namely
    $\Gamma_{L_0}(Z_jA^u)=\Gamma_{L_0}(A^jY_u)$. Injectivity of
    $\Gamma_{L_0}$, obtained from $L_0$-block injectivity, gives
    $Z_jA^u=A^jY_u$.
\end{proof}

\begin{theorem}[Inverting step for word compatibility]%
    \label{thm:right_factor_word_compatibility}
    \lean{MPSTensor.exists_right_factor_of_evalWord_compatibility}
    \leanok
    \uses{def:l_blk_injective}
    Assume $A$ is $L_0$-block-injective with $L_0 > 0$. If matrices $(Z_j)_j$
    satisfy $Z_jA^\sigma=A^jY_\sigma$ for every word $\sigma$ of length $K$,
    then there exists a matrix $X$ such that $Z_j=A^jX$ for all $j$.
\end{theorem}

\begin{proof}
    \leanok
    For $K = 1$, the hypothesis gives $Z_jA^i=A^jY_i$ for every letter $i$.
    Every blocked word of length $L_0$ is nonempty, so this identity extends to
    the coefficients of the blocked tensor $A^{[L_0]}$. Since $A^{[L_0]}$ is
    injective, the identity matrix is a linear combination of those blocked
    coefficients, and substituting this decomposition yields $Z_j=A^jX$ for a
    common matrix $X$. For larger $K$, strip the first letter: for each $i$, the
    matrices $Z_jA^i$ satisfy the same compatibility hypothesis with words of
    length $K-1$, so the induction hypothesis reduces the problem to the case
    $K=1$.
\end{proof}

\begin{theorem}[Growing-back step at the injectivity length]%
    \label{thm:ground_space_extend_right_block_injective}
    \lean{MPSTensor.groundSpace_extend_right_of_isNBlkInjective}
    \leanok
    \uses{def:in_left_ground, def:in_tail_ground, def:ground_space_parent}
    Assume $A$ is $L_0$-block-injective with $L_0 > 0$. If an element $\psi$
    on $K+L_0+1$ sites satisfies the left ground condition on the first
    $K+L_0$ sites and the suffix ground condition on every prefix of length
    $K$, then $\psi \in G_{K+L_0+1}(A)$.
\end{theorem}

\begin{proof}
    \leanok
    \uses{lem:left_tail_compatibility, thm:right_factor_word_compatibility}
    Lemma~\ref{lem:left_tail_compatibility} produces matrices $(Z_j)_j$ and
    $(Y_u)_u$ with $Z_jA^u=A^jY_u$ for every prefix word $u$ of length $K$.
    Theorem~\ref{thm:right_factor_word_compatibility} gives a common right
    factor $X$ such that $Z_j=A^jX$ for all $j$. Therefore, each last-site
    restriction of $\psi$ agrees with the corresponding last-site restriction
    of the ground-space vector determined by $X$. Since every configuration is
    obtained by adjoining its final letter to its initial $(K+L_0)$-tuple, the
    two states coincide, so $\psi$ lies in $G_{K+L_0+1}(A)$.
\end{proof}

\begin{definition}[Open-chain left and tail ground subspaces]
    \label{def:open_chain_ground_subspaces_es}
    \lean{MPSTensor.openChainLeftGroundSpaceES,
        MPSTensor.openChainTailGroundSpaceES,
        MPSTensor.openChainLeftGroundProjectionES,
        MPSTensor.openChainTailGroundProjectionES,
        MPSTensor.mem_openChainLeftGroundSpaceES_iff,
        MPSTensor.mem_openChainTailGroundSpaceES_iff}
    \leanok
    \uses{def:in_left_ground, def:in_tail_ground, def:ground_space_parent}
    In the Hilbert space of an open chain, let
    $\mathcal U_{K,L_0}(A)$ be the subspace satisfying the ground condition on
    the first $K+L_0$ sites, and let $\mathcal V_{K,L_0}(A)$ be the subspace
    satisfying the ground condition on the last $L_0+1$ sites. Equivalently,
    membership is tested by fixing respectively the final site or every
    prefix of length $K$ and applying the corresponding restriction map.
\end{definition}

\begin{theorem}[Open-chain left--tail ground-space intersection]
    \label{thm:open_chain_left_tail_ground_intersection}
    \lean{MPSTensor.openChainLeft_inf_tailGroundSpaceES}
    \leanok
    \uses{def:open_chain_ground_subspaces_es,
        thm:ground_space_extend_right_block_injective}
    If $A$ is $L_0$-block-injective and $L_0>0$, then
    \begin{align}
        \mathcal U_{K,L_0}(A)\cap\mathcal V_{K,L_0}(A)
        &=\mathcal G_{K+L_0+1}^{\mathrm{ES}}(A). \notag
    \end{align}
    In Nachtergaele's notation, with $n=K+L_0$ and $l=L_0$, this is the
    common range identity for $G_{\Lambda_n}$ and
    $G_{\Lambda_{n+1}\setminus\Lambda_{n-l}}$ inside $\Lambda_{n+1}$.
\end{theorem}

\begin{proof}
    \leanok
    The forward inclusion is the growing-back step at the injectivity length.
    Conversely, a vector in the longer ground space remains in the ground
    space after every left or tail restriction.
\end{proof}

\begin{theorem}[Open-chain iteration of the intersection property]%
    \label{thm:normal_open_chain_range_reduction}
    \lean{MPSTensor.contiguous_mem_groundSpace_of_isNBlkInjective}
    \leanok
    \uses{thm:ground_space_extend_right_block_injective, def:l_blk_injective,
        rem:contiguous_window_operators}
    Let $A$ be $L_0$-block-injective with $D \ge 1$ and $L_0 > 0$.
    If an $N$-site state satisfies the ground-space constraint on every
    non-wrapping contiguous interval of length $L_0+1$, for every fixed choice
    of the complementary sites, with $N \ge L_0+1$, then it lies in $G_N(A)$.
\end{theorem}

\begin{proof}
    \leanok
    \uses{thm:ground_space_extend_right_block_injective,
        def:tail_restrict_parent, rem:contiguous_window_operators}
    Induct on the extra length beyond $L_0+1$. The induction hypothesis gives
    the long left-window ground condition after deleting the last site, while
    Definition~\ref{def:tail_restrict_parent} and the contiguous-interval
    restriction identities above identify every fixed-prefix suffix with one
    of the assumed length-$L_0+1$ intervals.
    Theorem~\ref{thm:ground_space_extend_right_block_injective} then grows back
    the rightmost site.
\end{proof}

\begin{lemma}[Open-chain iteration for a sequence of subspaces]%
    \label{lem:open_chain_iteration_subspace_sequence}
    \lean{MPSTensor.contiguous_mem_of_restriction_intersection_submodules}
    \leanok
    \uses{rem:contiguous_window_operators}
    Let $S_m\subseteq(\C^d)^{\otimes m}$ be subspaces, and let $0<L\le N$.
    Suppose that every non-wrapping length-$L$ interval of $\psi$, for every
    fixed choice of the complementary sites, lies in $S_L$. Suppose also that,
    for every $m\ge L$,
    $(\C^d\otimes S_m)\cap (S_m\otimes\C^d) = S_{m+1}$.
    Then $\psi\in S_N$.
\end{lemma}

\begin{proof}
    \leanok
    \uses{rem:contiguous_window_operators}
    Induct on the length $m$ of a non-wrapping interval. The case $m=L$ is the
    assumed local condition. If all length-$m$ intervals lie in $S_m$, then the
    two one-site restrictions of a length-$(m+1)$ interval are adjacent
    length-$m$ intervals. Hence the length-$(m+1)$ interval lies in
    $(\C^d\otimes S_m)\cap(S_m\otimes\C^d)$, and the displayed identity places
    it in $S_{m+1}$. Taking $m=N$ and the interval starting at the first site
    gives $\psi\in S_N$.
\end{proof}

\begin{lemma}[Periodic constraints and propagated subspaces]%
    \label{lem:periodic_constraints_propagated_subspace_sequence}
    \lean{MPSTensor.chainGroundSpace_le_of_local_le_restriction_intersection_submodules}
    \leanok
    \uses{lem:open_chain_iteration_subspace_sequence, rem:cyclic_window_operators}
    Let $S_m\subseteq(\C^d)^{\otimes m}$ be subspaces, and let $0<L\le N$.
    Suppose that $G_L(A)\subseteq S_L$. Suppose also that, for every $m\ge L$,
    $(\C^d\otimes S_m)\cap (S_m\otimes\C^d) = S_{m+1}$.
    Then the periodic local constraints imply
    $\mathcal G_{N,L}(A)\subseteq S_N$.
\end{lemma}

\begin{proof}
    \leanok
    \uses{lem:open_chain_iteration_subspace_sequence, rem:cyclic_window_operators}
    If $\psi\in\mathcal G_{N,L}(A)$, then every length-$L$ cyclic interval of
    $\psi$ lies in $G_L(A)$, and hence in $S_L$. A non-wrapping interval is the
    cyclic interval starting at the same site. Therefore, all non-wrapping
    length-$L$ intervals lie in $S_L$, and
    Lemma~\ref{lem:open_chain_iteration_subspace_sequence} gives
    $\psi\in S_N$.
\end{proof}

\begin{lemma}[Cyclic-window range antitonicity]%
    \label{lem:cyclic_window_range_antitonicity}
    \lean{MPSTensor.chainGroundSpace_le_chainGroundSpace_succ,
        MPSTensor.chainGroundSpace_le_chainGroundSpace_of_le}
    \leanok
    \uses{rem:cyclic_window_operators}
    On a nonempty periodic chain, longer cyclic-window constraints imply all
    shorter cyclic-window constraints: if $L'\le L\le N$, then
    $\mathcal G_{N,L}(A)\subseteq\mathcal G_{N,L'}(A)$.
\end{lemma}

\begin{proof}
    \leanok
    \uses{lem:ground_space_in_left_ground}
    For every cyclic window vector $v\in G_{L+1}(A)$ and every last letter $b$,
    one has $v|_{\mathrm{last}=b}\in G_L(A)$. Thus
    $\mathcal G_{N,L+1}(A)\subseteq\mathcal G_{N,L}(A)$, and iteration gives
    the general antitonicity.
\end{proof}

\begin{theorem}[Open-chain containment from cyclic local constraints]%
    \label{thm:cyclic_normal_open_chain_range_reduction}
    \lean{MPSTensor.chainGroundSpace_le_groundSpace_of_isNBlkInjective}
    \leanok
    \uses{lem:cyclic_window_range_antitonicity,
        thm:normal_open_chain_range_reduction}
    Let $A$ be $L_0$-block-injective with $L_0>0$ and $D \ge 1$.
    If an $N$-site periodic-chain state satisfies all cyclic ground-space
    constraints at some range $L_0+1\le L\le N$, then it lies in the
    open-chain ground space $G_N(A)$.
\end{theorem}

\begin{proof}
    \leanok
    \uses{lem:cyclic_window_range_antitonicity,
        thm:normal_open_chain_range_reduction}
    Lemma~\ref{lem:cyclic_window_range_antitonicity} reduces the cyclic
    hypotheses to range $L_0+1$. On every non-wrapping interval, the cyclic and
    contiguous restriction maps agree, so
    Theorem~\ref{thm:normal_open_chain_range_reduction} applies.
\end{proof}
